\chapter{Lecture des données de marché}\label{ch:m1:reading-market-data}

Avant midi, une bourse a envoyé à une même firme des dizaines de millions de messages portant sur
des milliers d'instruments. Chacun annonce qu'un ordre a été ajouté, réduit,
exécuté ou retiré, en quelques dizaines d'octets. Aucun ne dit quel est le meilleur cours acheteur~:
c'est au destinataire de le déduire, en appliquant chaque message, dans l'ordre, à
un carnet d'ordres qu'il maintient lui-même. Si un message se perd, le carnet est faux
jusqu'à la fin de la journée, en silence. Si un compte rendu de transaction comporte une faute de frappe, un cours
de 10.01 pour une action valant 100, le plus bas de la journée, sa volatilité, son
prix moyen pondéré par les volumes et tout signal calculé à partir de ces données sont faux,
tout aussi en silence. Chaque résultat de la suite de cette série repose sur un tableau de
cours~; ce chapitre porte sur la façon dont ce tableau est construit et sur les défauts
dont il hérite.

\section{Trois niveaux de données}

\begin{definition}[Niveau 1, niveau 2, niveau 3]\label{def:m1:reading-market-data:levels}
Les données de \emph{niveau 1}\index{niveau 1} sont le meilleur cours acheteur et le meilleur cours vendeur, avec leurs
quantités, et la dernière transaction. Les données de \emph{niveau 2}\index{niveau 2} sont le carnet
d'ordres agrégé par prix~: pour chaque palier de prix, la quantité totale (et parfois
le nombre d'ordres), sur une certaine profondeur ou en totalité~: le flux \omterm{def:m1:matching-algorithms-and-implied-spreads:mbo}{market-by-price} de
la \cref{def:m1:matching-algorithms-and-implied-spreads:mbo}. Les données de \emph{niveau
3}\index{niveau 3} sont la totalité des ordres, pris individuellement, avec leur identifiant et
chacun de leurs événements au cours de leur vie~: le flux \omterm{def:m1:matching-algorithms-and-implied-spreads:mbo}{market-by-order}, à partir duquel les deux autres
peuvent être calculés, et seul lui donne les rangs dans la file d'attente.
\end{definition}

\begin{omfigure}
\begin{tikzpicture}[font=\small,
  col/.style={draw=omDef,rounded corners=2pt,minimum width=3.9cm,minimum height=3.3cm,align=left,fill=omDef!4,font=\footnotesize\ttfamily,inner sep=5pt}]
\node[col] (l3) at (0,0) {réf. côté  quantité  prix\\1001 A \ \ \ 300 \ 99.99\\1002 A \ \ \ 200 \ 99.99\\1004 A \ \ \ 100 \ 99.98\\1003 V \ \ \ 500 100.01\\1005 V \ \ \ 100 100.02};
\node[col] (l2) at (4.6,0) {côté  quantité  prix  nb\\A \ \ \ \ 500 \ 99.99 \ 2\\A \ \ \ \ 100 \ 99.98 \ 1\\V \ \ \ \ 500 100.01 \ 1\\V \ \ \ \ 100 100.02 \ 1};
\node[col] (l1) at (9.2,0) {bid \ 99.99 x 500\\ask 100.01 x 500\\dernier 100.01 x 100};
\node[font=\small\bfseries,above=2pt of l3] {Niveau~3~: ordres};
\node[font=\small\bfseries,above=2pt of l2] {Niveau~2~: prix};
\node[font=\small\bfseries,above=2pt of l1] {Niveau~1~: meilleurs cours};
\draw[-{Stealth[length=2mm]},thick] (l3.east) -- (l2.west);
\draw[-{Stealth[length=2mm]},thick] (l2.east) -- (l1.west);
\end{tikzpicture}

\omcaption{Un même carnet à trois niveaux de détail. Chaque niveau peut être calculé
à partir de celui situé à sa gauche, et jamais l'inverse~: les barres d'une seconde d'un fournisseur
sont des données de \omterm{def:m1:reading-market-data:levels}{niveau 1} échantillonnées~; elles ne peuvent pas être reconstituées en un
carnet d'ordres.}\label{fig:m1:reading-market-data:levels}
\end{omfigure}

\begin{definition}[Numéro de séquence]\label{def:m1:reading-market-data:seq}
Un \emph{numéro de séquence}\index{numéro de séquence} est un compteur que le
diffuseur incrémente à chaque message sur un canal. Un récepteur qui constate un
saut dans la numérotation a perdu des messages et doit les récupérer, auprès d'un service
de retransmission ou d'un instantané, avant que son carnet puisse de nouveau être jugé fiable.
\end{definition}

Un flux de \omterm{def:m1:reading-market-data:levels}{niveau 3} est un flux de \emph{différences}. L'état se trouve du
côté du récepteur, ce qui rend le protocole compact et fait de chaque récepteur
responsable de l'exactitude des données. Les \omterm{def:m1:us-equity-market-structure:sip}{flux directs} sont publiés sur deux
chemins réseau redondants~; le récepteur retient, pour chaque \omterm{def:m1:reading-market-data:seq}{numéro de séquence}, la copie
qui arrive en premier~; il signale le cas où aucune des deux n'arrive.

\begin{dated}{2026-09}{Un format de transmission}\label{dat:m1:reading-market-data:itch}
Le flux actions ordre par ordre dont la spécification version 5.0 sert de
modèle à ce chapitre encode tous les entiers en big-endian~; les prix sont des entiers
avec quatre décimales implicites~; les horodatages en nanosecondes depuis minuit
sur six octets~; et les instruments par un code \emph{locate} de deux octets attribué pour
la journée, à la même position dans chaque message. Un message d'ajout d'ordre compte 36
octets~: type, locate, numéro de suivi, horodatage, référence de l'ordre sur huit
octets, côté, nombre d'actions, un symbole de huit caractères et le prix.
Une exécution compte 31 octets et ne renvoie à l'ordre que par sa référence, comme le font
l'annulation (23 octets) et la suppression (19 octets)~: le prix d'une transaction ne figure pas
dans le message d'exécution~: il se trouve là où le récepteur a stocké l'ordre.
\end{dated}

\section{Transactions, cotations et codes de condition}

\begin{definition}[Condition de transaction]\label{def:m1:reading-market-data:condition}
Une \emph{condition de transaction}\index{condition de transaction} (ou condition de vente) est un code
rattaché à un compte rendu de transaction qui indique la nature de celle-ci et la manière dont elle
peut être utilisée~: une transaction ordinaire~; une transaction d'ouverture ou de clôture en enchère~; un \omterm{def:m1:us-equity-market-structure:lots}{lot
irrégulier}~; une transaction déclarée en retard ou hors séquence~; une transaction à un prix
moyen ou dérivé~; une correction ou une annulation d'un compte rendu antérieur. Les
consolidated tapes et chaque bourse définissent les conditions qui mettent à jour le dernier
cours, le plus haut et le plus bas, et le volume.
\end{definition}

Un tape n'est pas une liste de cours auxquels on aurait pu négocier. Un
compte rendu hors séquence porte un cours datant de quelques minutes. Une
transaction à prix moyen ne correspond à aucune cotation qui ait jamais existé. Une transaction annulée n'a pas eu lieu~;
son annulation arrive sous forme d'un message ultérieur qu'un chargeur négligent ignore.
Une transaction d'enchère est un seul cours pour une large part du volume de la journée. Les transactions hors bourse
(décrites au \cref{ch:m1:retail-flow-and-wholesaling}) sont déclarées après
coup, avec quelques secondes de retard, et un horodatage pouvant être celui du déclarant.
Avant de calculer la moindre statistique à partir des transactions, une règle doit préciser quelles
conditions sont retenues.

\begin{method}[Nettoyage d'un tape]\label{met:m1:reading-market-data:clean}
\begin{enumerate}
\item Appliquer les corrections et les annulations en les rapprochant des
comptes rendus d'origine.
\item Conserver ou écarter explicitement chaque condition, selon l'usage~: tout le volume exécuté
pour le volume total~; seulement les transactions régulières et en séquence pour le
dernier cours, un plus haut ou un plus bas~; les enchères à part.
\item Signaler les transactions très éloignées de la cotation en vigueur ou d'une médiane glissante~;
les examiner avant de les supprimer~: un flash crash est une donnée, une erreur de décimale n'en est pas une.
\item Consigner ce qui a été supprimé et pourquoi. La règle fait partie du résultat.
\end{enumerate}
\end{method}

\begin{omfigure}
\begin{tikzpicture}
\begin{axis}[width=10.5cm,height=4.8cm,
  xlabel={numéro de transaction},ylabel={prix},
  tick label style={font=\small},label style={font=\small},
  xmin=0,xmax=16,ymin=0,ymax=115,grid=major,grid style={gray!25},
  scatter/classes={1={mark=*,omDef,mark size=1.8pt},0={mark=x,omThm,mark size=3.5pt,very thick}},
  legend columns=2,legend style={font=\footnotesize,at={(0.5,-0.32)},anchor=north,draw=none,fill=none,/tikz/every even column/.append style={column sep=8pt}},
  table/col sep=comma]
\addplot[scatter,only marks,scatter src=explicit symbolic] table[x=seq,y=price,meta=kept]{figdata/markets-1/28-reading-market-data/badtick.csv};
\legend{conservées,écartées par le filtre}
\end{axis}
\end{tikzpicture}

\omcaption{Quinze transactions déclarées sur un même titre. Le filtre automatique en écarte
deux, l'impression à 10.01 et une transaction annulée d'apparence tout à fait normale~;
la déclaration tardive à 99.80 passe à travers. Avec les quinze, le prix moyen pondéré
par les volumes vaut 98.40~; la valeur réelle est proche de 100. Données~: le
problème du week-end.}\label{fig:m1:reading-market-data:badtick}
\end{omfigure}

\section{Horodatages}

\begin{definition}[Horodatages de la bourse et de réception]\label{def:m1:reading-market-data:ts}
L'\emph{horodatage de la bourse}\index{horodatage de la bourse} est écrit par le
moteur d'appariement ou la passerelle de l'éditeur lorsque l'événement se produit. Le
\emph{horodatage de réception}\index{horodatage de réception} est écrit par le
système de capture du destinataire lorsque le message arrive. Le premier ordonne les événements
au sein d'une même plateforme~; seul le second indique ce que le destinataire pouvait savoir,
et à quel moment.
\end{definition}

Un backtest qui utilise les \omterm{def:m1:reading-market-data:ts}{horodatages de la bourse} suppose que la firme a pris connaissance de chaque
événement à l'instant où il s'est produit. L'écart, de quelques dizaines de microsecondes dans
un même bâtiment et de quelques millisecondes entre villes, constitue à lui seul
l'avantage de certaines stratégies (livre~10). Les horloges de deux plateformes ne sont pas
la même horloge~: comparer leurs horodatages pour décider quelle cotation a bougé en premier
n'est valable qu'à la précision de leur synchronisation.

\begin{dated}{2026-09}{Précision requise des horloges}\label{dat:m1:reading-market-data:clocks}
\emph{Union européenne.} La norme technique de réglementation sur la synchronisation des horloges
impose aux membres d'une plateforme de négociation qui recourent à des techniques
algorithmiques à haute fréquence de maintenir leurs horloges de gestion à moins de 100
microsecondes de l'UTC et d'horodater avec une granularité d'une microseconde ou
mieux. \emph{États-Unis.} Le plan du consolidated audit trail exige que les bourses
maintiennent leurs horloges à moins de 100 microsecondes de l'heure de l'institut national
de \omterm{def:m1:reading-market-data:norm}{normalisation}, et que les \omterm{def:m1:the-sell-side:sell-side}{broker-dealers} le fassent à moins de 50 millisecondes (une
seconde pour les horloges qui ne servent qu'aux événements manuels).
\end{dated}

\begin{omfigure}
\begin{tikzpicture}
\begin{axis}[width=10cm,height=4.6cm,
  xlabel={gigue moyenne de latence (microsecondes)},ylabel={hors d'ordre (\%)},
  tick label style={font=\small},label style={font=\small},
  xmin=0,xmax=5200,ymin=0,ymax=36,grid=major,grid style={gray!25},table/col sep=comma]
\addplot[omDef,very thick,mark=*,mark size=1.4pt] table[x=jitter_us,y=inversions_pct]{figdata/markets-1/28-reading-market-data/reorder.csv};
\end{axis}
\end{tikzpicture}

\omcaption{L'échantillon de flux du chapitre, rejoué avec un \omterm{def:m1:reading-market-data:ts}{horodatage de réception} égal
à l'\omterm{def:m1:reading-market-data:ts}{horodatage de la bourse}, plus 200 microsecondes, plus une gigue exponentielle. Trier
par \omterm{def:m1:reading-market-data:ts}{horodatage de réception} désordonne 1~\% des messages pour une gigue de 100 microsecondes et
un tiers pour 5 millisecondes. Au sein d'un même flux, ordonner par \omterm{def:m1:reading-market-data:seq}{numéro de séquence}, jamais
par l'une ou l'autre horloge. Données~: le tutoriel.}\label{fig:m1:reading-market-data:reorder}
\end{omfigure}

\section{Symbologie et données de référence}

\begin{definition}[Symbologie et données de référence]\label{def:m1:reading-market-data:refdata}
La \emph{symbologie}\index{symbologie} est l'ensemble des identifiants sous lesquels un
même instrument est connu~: symbole boursier de cotation, codes des fournisseurs, numéros nationaux et
internationaux de valeurs, code numérique du flux pour la journée.
Les \emph{données de référence}\index{données de référence} regroupent tout ce qui concerne
l'instrument et qui n'est pas un prix~: identifiants et leur historique, plateforme de cotation,
devise, pas de cotation, taille de lot, multiplicateur, échéance, opérations sur titres,
calendrier de négociation.
\end{definition}

Un symbole boursier n'est pas un identifiant. Les symboles sont réutilisés après des radiations, modifiés
lors de fusions, et diffèrent selon la plateforme et le fournisseur~; le code locate d'un flux n'est valable
pour une journée. Les tables de recherche doivent être indexées par un identifiant interne
permanent, chaque symbole externe étant rattaché à celui-ci par plage de dates. L'erreur
silencieuse la plus fréquente en recherche sur actions, après le look-ahead, est une
série temporelle qui assemble deux entreprises différentes à la date où l'une a repris le
symbole de l'autre. Les constructions des chapitres~\ref{ch:m1:shares-corporate-actions-indices},
\ref{ch:m1:futures-contracts-and-exchanges} et
\ref{ch:m1:option-contracts-and-exchanges} constituent les \omterm{def:m1:reading-market-data:refdata}{données de référence} de ce chapitre.

\section{Du flux brut à la table de recherche}

\begin{definition}[Normalisation]\label{def:m1:reading-market-data:norm}
La \emph{normalisation}\index{normalisation} est la conversion de chaque plateforme
de ses messages natifs en une représentation interne unique~: types de messages communs,
prix entiers dans une unité commune, champ d'horloge commun, identifiants permanents
des instruments, indicateurs explicites des conditions.
\end{definition}

\begin{omfigure}
\begin{tikzpicture}
\begin{axis}[width=11cm,height=5cm,
  xlabel={secondes depuis le premier message},ylabel={prix},
  tick label style={font=\small},label style={font=\small},
  ymin=99.945,ymax=100.065,grid=major,grid style={gray!25},
  yticklabel style={/pgf/number format/fixed,/pgf/number format/fixed zerofill,/pgf/number format/precision=2},
  legend columns=3,legend style={font=\footnotesize,at={(0.5,-0.3)},anchor=north,draw=none,fill=none,/tikz/every even column/.append style={column sep=8pt}},
  table/col sep=comma]
\addplot[omDef,thick,const plot] table[x=seconds,y=bid]{figdata/markets-1/28-reading-market-data/l1.csv};
\addplot[omThm,thick,const plot] table[x=seconds,y=ask]{figdata/markets-1/28-reading-market-data/l1.csv};
\addplot[only marks,mark=*,mark size=1.2pt,black] table[x=seconds,y=price]{figdata/markets-1/28-reading-market-data/tape.csv};
\legend{meilleur bid,meilleur ask,transactions}
\end{axis}
\end{tikzpicture}

\omcaption{\omterm{def:m1:reading-market-data:levels}{Niveau 1} et le tape, reconstruits à partir de 640 messages de l'échantillon
binaire du chapitre, par le décodeur et le carnet de la construction. Chaque exécution d'un ordre affiché
s'imprime au bid ou à l'ask, ce qui permet de connaître son agresseur~;
les impressions intermédiaires sont des transactions contre des ordres cachés. Données~: la
construction du chapitre.}\label{fig:m1:reading-market-data:l1}
\end{omfigure}

Le pipeline comporte quatre étapes, et la recherche doit savoir de laquelle
provient sa table. \emph{Capture}~: chaque paquet, avec un \omterm{def:m1:reading-market-data:ts}{horodatage de réception} matériel,
écrit sans modification. \emph{Décodage et \omterm{def:m1:reading-market-data:norm}{normalisation}}~: la construction ci-dessous.
\emph{Construction du carnet}~: appliquer les messages dans l'ordre de séquence, détecter les trous, émettre le \omterm{def:m1:reading-market-data:levels}{niveau 2}
et le \omterm{def:m1:reading-market-data:levels}{niveau 1} comme flux dérivés. \emph{Tables}~: transactions avec conditions et
côté agresseur~; cotations~; barres, qui constituent un résumé avec perte choisi dans un
but précis. Chaque étape doit pouvoir être rejouée à partir de la précédente, bit pour bit~:
lorsqu'un résultat paraît faux, la question «~est-ce la donnée~?~» doit avoir une
réponse.

\section{Tutoriel~: un flux binaire alimentant un carnet et un tape}\label{tut:m1:reading-market-data:decode}

\begin{tutorial}
\textbf{Objectif.} Décoder un flux binaire ordre par ordre encadré, tenir à jour le carnet,
produire le \omterm{def:m1:reading-market-data:levels}{niveau 1} et un tape avec le côté agresseur, et mesurer l'effet de la gigue de
l'\omterm{def:m1:reading-market-data:ts}{horodatage de réception} sur l'ordre des messages. \textbf{État final~:} les trois figures de
données de ce chapitre, à partir de \texttt{code/firm/feed/data/sample.itch} (2\,020 messages,
65\,842 octets).
\begin{tutsteps}
\item \textbf{Encadrement et en-tête.} Préfixe de longueur, type, puis les champs que
tous les messages partagent. Une longueur incompatible avec le type est une erreur, non
une indication.
\omcode{code/firm/feed/firm_feed.py}{58}{74}{Python~: encadrer, contrôler, lire l'en-tête commun.}\label{lst:m1:reading-market-data:py}
\item \textbf{Le même code en C++}, sur un span d'octets, sans copies.
\omcode{code/firm/feed/cpp/firm_feed.hpp}{47}{59}{C++~: l'encadrement et l'en-tête de la boucle de décodage.}\label{lst:m1:reading-market-data:cpp}
\item \textbf{Et en Rust}, sous la forme d'un itérateur qui produit un résultat par message
et s'arrête après la première erreur.
\omcode{code/firm/feed/rust/src/lib.rs}{53}{79}{Rust~: une étape du décodeur.}\label{lst:m1:reading-market-data:rs}
\item \textbf{Le carnet.} Une exécution ne porte aucun prix~: il faut le retrouver.
\omcode{code/firm/feed/firm_feed.py}{146}{168}{Application d'un message au carnet et au tape.}\label{lst:m1:reading-market-data:apply}
\end{tutsteps}
\textbf{Pour aller plus loin.} Retirer un message d'ajout d'ordre de l'échantillon et
observer les erreurs que le carnet enregistre, ainsi que le délai avant qu'une exécution
ne fasse référence à l'ordre manquant. Ajouter ensuite un message de remplacement (annulation
et ajout avec une nouvelle référence, avec perte de priorité) dans les trois décodeurs.
\end{tutorial}

\section{Construction~: le normaliseur de flux}\label{bld:m1:reading-market-data:feed}

\begin{build}
\bfield{Rôle} Le point d'entrée de toutes les données de marché dans la
mini-firme~: la bourse simulée du livre~10 publie ce format, les
stratégies du livre~11 consomment le carnet, et les tables de recherche du livre~7 sont
écrites à partir du tape.
\bfield{Interface} Python: \texttt{encode(msg)}, \texttt{decode(buffer)},
\texttt{Book.apply(msg)}, \texttt{Book.best(locate)}, \texttt{Book.depth(locate,
side, n)}, \texttt{Book.trades}, \texttt{Book.errors}, \texttt{summary(book,
locate)}.
C++20 (en-tête seul)~: \texttt{firm::feed::decode(span, callback)},
\texttt{Book}, \texttt{Summary}. Rust~: \texttt{decode(\&[u8])} renvoyant un
itérateur de résultats, \texttt{Book}, \texttt{Summary}.
\bfield{Règles} Entiers en gros-boutiste, prix en dix-millièmes, horodatages de
48~bits en nanosecondes, encadrement par préfixe de longueur sur deux octets. Une trame
tronquée ou une longueur incompatible avec le type arrête le décodage par une erreur. Le
carnet ne devine jamais~: une référence d'ordre inconnue, une sur-exécution, un ajout en double
ou un horodatage qui recule est enregistré et compté. Les trois implémentations doivent
produire le même résumé sur l'échantillon commun.
\bfield{Tests d'acceptation} \texttt{code/firm/feed/tests/} (aller-retour, données
aléatoires, un carnet construit à la main, erreurs enregistrées, reproductibilité de
l'échantillon)~; \texttt{cpp/firm\_feed\_test.cpp} et les tests de la crate Rust (le
résumé de l'échantillon commun, la troncature, une mauvaise longueur, un ordre inconnu).
\bfield{Pour aller plus loin} \omterm{def:m1:reading-market-data:seq}{Numéros de séquence} avec détection des trous et arbitrage
des lignes A/B~; un message d'instantané pour la reprise~; un benchmark du nombre de messages
par seconde dans chaque langage (le livre~13 part de là).
\end{build}

\begin{omsources}
\item Nasdaq, \emph{Nasdaq TotalView-ITCH 5.0} spécification (révision du 28
avril 2023).
\item Commission Delegated Regulation (EU) 2017/574 (synchronisation des horloges,
«~RTS 25~»).
\item CAT NMS Plan, questions fréquemment posées sur la synchronisation des horloges
(R1)~; FINRA Rule 6820.
\end{omsources}

\section{Exercices}

\begin{exercise}[$\star$]\label{exo:m1:reading-market-data:1}
À partir de la vue de \omterm{def:m1:reading-market-data:levels}{niveau 3} de la \cref{fig:m1:reading-market-data:levels}, une exécution
de 300 actions arrive pour l'ordre 1001, puis une annulation de 100 actions pour l'ordre
1002. Donner les niveaux 2 et 1 après ces événements. À quel prix la transaction a-t-elle
été publiée, et qui était l'agresseur~?
\end{exercise}

\begin{exercise}[$\star$]\label{exo:m1:reading-market-data:2}
Un champ de prix contient les quatre octets \texttt{00 0F 41 DC} (hexadécimal),
en gros-boutiste, avec quatre décimales implicites. Donner le prix. Que lirait un
lecteur en petit-boutiste~?
\end{exercise}

\begin{exercise}[$\star$]\label{exo:m1:reading-market-data:3}
L'échantillon du chapitre comporte 1\,102 ajouts, 311 exécutions, 114 annulations, 473 suppressions
et 20 transactions cachées. Vérifier la taille du fichier à partir des longueurs des messages
et de l'encadrement.
\end{exercise}

\begin{exercise}[$\star\star$]\label{exo:m1:reading-market-data:4}
Un horodatage de six octets compte les nanosecondes depuis minuit. Quelle est la plus grande
heure de la journée qu'il peut contenir~? Pourquoi quatre octets ne suffisent-ils pas, et que
gagne-t-on par rapport à huit~?
\end{exercise}

\begin{exercise}[$\star\star$]\label{exo:m1:reading-market-data:5}
La plateforme A horodate une modification de cotation à 10:00:00.000\,150 et la plateforme B
à 10:00:00.000\,190. Les horloges des deux plateformes sont à moins de 100 microsecondes de
l'UTC. Peut-on dire laquelle a bougé la première~? De quoi aurait-on besoin~?
\end{exercise}

\begin{exercise}[$\star\star$]\label{exo:m1:reading-market-data:6}
Un récepteur voit les \omterm{def:m1:reading-market-data:seq}{numéros de séquence} 5\,001, 5\,002, 5\,004 sur la ligne A et 5\,001,
5\,003, 5\,004 sur la ligne B. Que fait-il~? Maintenant, les deux lignes sautent 5\,003~:
que faut-il faire avant de réutiliser le carnet, et que peut faire la stratégie
entre-temps~?
\end{exercise}

\begin{exercise}[$\star\star\star$]\label{exo:m1:reading-market-data:7}
\emph{Programmation.} Avec \texttt{receive\_times} et \texttt{inversions} sur l'échantillon,
reproduire les pourcentages de messages hors ordre pour des gigues de 100 et 1\,000
microsecondes. Diviser ensuite le débit de messages par deux (garder un message sur deux) et
expliquer la variation.
\end{exercise}

\begin{exercise}[$\star\star\star$]\label{exo:m1:reading-market-data:8}
\emph{Chasse à l'erreur.} Un chercheur télécharge des barres d'une minute pour 3\,000
tickers sur quinze ans, les indexe par ticker, calcule le plus bas quotidien de chaque action
comme le minimum des plus bas des barres, et constate qu'acheter les actions qui ont chuté de
plus de 50~\% en cours de séance et remonté à la clôture est extrêmement rentable. Nommer
trois défauts de données, vus dans ce chapitre, qui pourraient produire ce résultat.
\end{exercise}

\section{Problème~: le mauvais tick}

\begin{problem}[{Problème du week-end --- quinze impressions et un indice de référence}]\label{pb:m1:reading-market-data:1}
Un desk doit déclarer le prix moyen pondéré par les volumes d'une action entre 09:30:00
et 09:32:00, en tant qu'indice de référence d'un ordre client. Le ruban (séquence, heure,
prix, actions, condition~; les conditions sont les étiquettes de ce livre)~:

\begin{center}\footnotesize
\begin{tabular}{rlrrl@{\qquad}rlrrl}
\toprule
1 & 09:30:00 & 100.02 & 12\,000 & ouverture & 9 & 09:30:52 & 100.09 & 300 & annulée \\
2 & 09:30:04 & 100.05 & 300 & & 10 & 09:31:03 & 100.11 & 600 & \\
3 & 09:30:09 & 100.04 & 50 & \omterm{def:m1:us-equity-market-structure:lots}{lot irrégulier} & 11 & 09:31:10 & 99.80 & 5\,000 & tardive \\
4 & 09:30:15 & 100.07 & 500 & & 12 & 09:31:18 & 100.12 & 200 & \\
5 & 09:30:21 & 100.06 & 200 & & 13 & 09:31:25 & 100.13 & 700 & \\
6 & 09:30:30 & 10.01 & 400 & & 14 & 09:31:40 & 100.12 & 100 & \\
7 & 09:30:31 & 100.08 & 400 & & 15 & 09:31:55 & 100.15 & 900 & \\
8 & 09:30:40 & 100.10 & 1\,000 & & & & & & \\
\bottomrule
\end{tabular}
\end{center}

\medskip\noindent\textbf{Partie I --- Le chiffre naïf.}
\begin{enumerate}
\item Donner le total des actions et le prix moyen pondéré par les volumes de l'ensemble des
quinze impressions.
\item Donner le plus bas de la période à partir de toutes les impressions.
\item De combien de points de base la moyenne naïve diffère-t-elle de 100.05~?
\item Que produirait une estimation de la volatilité à partir de ces impressions~?
\end{enumerate}

\medskip\noindent\textbf{Partie II --- Impression par impression.}
\begin{enumerate}[resume]
\item Impression 6~: de quoi s'agit-il, comment un filtre la détecterait-il, et comment
le confirmer~?
\item Impression 9~: pourquoi un filtre sur les valeurs aberrantes ne la détecte-t-il pas, et qu'est-ce qui le fait~?
\item L'impression 11 a été exécutée à 09:24 et déclarée à 09:31. Doit-elle figurer dans
l'indice de référence d'un ordre commencé à 09:30~?
\item L'impression 3 est un \omterm{def:m1:us-equity-market-structure:lots}{lot irrégulier}. Dedans ou dehors~?
\item L'impression 1 est l'enchère d'ouverture. L'ordre du client a été saisi à
09:30:02. Dedans ou dehors~?
\end{enumerate}

\medskip\noindent\textbf{Partie III --- Les chiffres corrigés.}
\begin{enumerate}[resume]
\item En excluant les impressions erronées, annulées et tardives, donner le nombre
d'impressions exclues, les actions et le prix moyen.
\item En excluant également l'enchère, donner le prix moyen de la
négociation en continu.
\item Donner le plus bas et le plus haut corrigés.
\item L'ordre du client a acheté 2\,000 actions à un prix moyen de 100.09. Donner son
glissement en points de base par rapport à chacun des deux indices corrigés.
\item Lequel des deux indices est le bon~?
\end{enumerate}

\medskip\noindent\textbf{Partie IV --- Jugement.}
\begin{enumerate}[resume]
\item La bourse annule plus tard la transaction~6. Comment cela parvient-il à votre base de
données, et que doit faire votre chargeur~?
\item Pourquoi faut-il fixer la règle de nettoyage avant d'examiner
le résultat de l'ordre~?
\item Les barres «~nettoyées~» d'un fournisseur excluent déjà certaines conditions. Que
devez-vous obtenir du fournisseur~?
\item À quelle étape du pipeline en quatre étapes chacun des défauts de la partie II est-il
traité~?
\item Énoncer le \emph{résultat nommé}~: le nombre d'impressions à exclure et le prix moyen
pondéré par les volumes corrigé.
\item En une phrase~: qu'est-ce qu'un compte rendu de transaction~?
\end{enumerate}
\end{problem}

\section{Questions d'entretien}

\begin{interviewq}[$\star$ \iqroles{développeur, chercheur, trader}]\label{iq:m1:reading-market-data:1}
Quelle est la différence entre les données de marché de \omterm{def:m1:reading-market-data:levels}{niveau 1}, de \omterm{def:m1:reading-market-data:levels}{niveau 2} et de \omterm{def:m1:reading-market-data:levels}{niveau 3}~?
\end{interviewq}

\begin{interviewq}[$\star$ \iqroles{développeur}]\label{iq:m1:reading-market-data:2}
Comment construisez-vous un carnet d'ordres à partir d'un flux ordre par ordre~?
Quelles structures de données utilisez-vous~?
\end{interviewq}

\begin{interviewq}[$\star\star$ \iqroles{développeur, chercheur}]\label{iq:m1:reading-market-data:3}
Votre gestionnaire de flux détecte un trou dans les \omterm{def:m1:reading-market-data:seq}{numéros de séquence}. Que se passe-t-il ensuite~?
\end{interviewq}

\begin{interviewq}[$\star\star$ \iqroles{chercheur, ingénieur ML}]\label{iq:m1:reading-market-data:4}
Quel horodatage utilisez-vous dans un backtest, et que se passe-t-il si vous choisissez
l'autre~?
\end{interviewq}

\begin{interviewq}[$\star\star$ \iqroles{chercheur, ingénieur ML}]\label{iq:m1:reading-market-data:5}
On vous remet un fichier de transactions. Que vérifiez-vous avant de calculer quoi que ce soit~?
\end{interviewq}

\begin{interviewq}[$\star\star\star$ \iqroles{développeur}]\label{iq:m1:reading-market-data:6}
Concevoir le stockage des données de marché captées par une société de sorte que tout résultat de recherche
puisse être reproduit cinq ans plus tard.
\end{interviewq}
